
%\medskip
%
%\parbox{0.8\linewidth}{Hugo réalise un assemblage de carreaux représentant son héros préféré. 
%
%Pour cela il doit coller 22 carreaux violets, 2 blancs, 162 noirs et 110 verts.
%
%Tous les carreaux sont mélangés dans une boîte.
%
%Hugo choisit un carreau au hasard.
%
%On estime que tous les carreaux ont la même chance d'être choisis.

\medskip
\begin{enumerate}
\item %Quelle est la probabilité que Hugo choisisse un carreau vert ?
Il y a 110 carreaux verts sur un total de $22 + 2 + 162 + 110 = 296$ carreaux.

La probabilité de tirer un carreau vert est égale à $\dfrac{110}{296} = \dfrac{55}{148}$.
\item %Quelle est la probabilité que Hugo ne choisisse pas un carreau violet ?
La probabilité de choisir un carreau violet est $\dfrac{22}{296} = \dfrac{11}{148}$, donc la probabilité de ne pas choisir un carreau violet est $1 - \dfrac{11}{148} = \dfrac{148 - 11}{148} = \dfrac{137}{148}$.
\item %Quelle est la probabilité que le carreau choisi soit noir ou blanc?
La probabilité que le carreau choisi soit noir ou blanc est $\dfrac{162 + 2}{296} = \dfrac{164}{296} = \dfrac{41}{74}$.
\item %En une journée Hugo a collé 75\,\% des carreaux. Combien de carreaux cela représente-t-il ?
On a $\dfrac{75}{100} \times 296 = \dfrac{\np{22200}}{100} = 222$.

Hugo a collé 222 carreaux en une journée.
\end{enumerate}
%} 
%\hfill \parbox{0.18\linewidth}{\includegraphics[width=3cm]{Heros}}

\vspace{0.5cm}

